\chapter{Tail-Dominated Outliers and Activation Range Control in Low-Bit Training}

\section{Introduction}

When training large Transformer models under low-bit quantization (e.g., FP8, FP4, INT4), one of the most persistent and difficult challenges is the presence of activation outliers. These outliers are individual coordinates within token embeddings whose magnitudes are significantly larger than typical values in the same vector. Although such extreme values may appear harmless from a floating-point perspective, they become highly problematic under finite precision. In uniform quantization schemes, the quantization step size is determined by the maximum absolute value within the vector. As a consequence, a small number of extreme coordinates can inflate the dynamic range, forcing the quantizer to allocate large bins and thereby coarsening the resolution available for the majority of coordinates. This effect disproportionately degrades small-magnitude but semantically meaningful signals.

To design a principled solution, we must first understand the structural origin of these activation outliers. In particular, we must answer two questions. First, do these outliers originate from the dominant, low-rank structure of the representation? Or do they arise from the long-tail components? Second, how does this distinction affect quantization robustness? Addressing these questions requires a careful examination of the spike–tail decomposition of Transformer representations after mean removal.

\section{Spike–Tail Decomposition After Mean Removal}

Let $x \in \mathbb{R}^d$ denote a token embedding after mean removal. Across tokens in a training batch or dataset, we may compute the covariance matrix

\[
\Sigma = \mathbb{E}[x x^\top].
\]

Empirical studies of Transformer activations consistently reveal that $\Sigma$ is highly anisotropic. Specifically, its eigenvalue spectrum exhibits a small number of dominant eigenvalues followed by a long decaying tail. Writing the eigendecomposition as

\[
\Sigma = U \Lambda U^\top,
\]

where $U$ contains orthonormal eigenvectors and $\Lambda$ is diagonal with eigenvalues ordered $\lambda_1 \ge \lambda_2 \ge \dots \ge \lambda_d$, we observe that $\lambda_1, \dots, \lambda_k$ are significantly larger than the rest for a modest $k \ll d$.

This motivates the decomposition of each token embedding into two orthogonal components:

\[
x = x_{\text{spike}} + x_{\text{tail}},
\]

where

\[
x_{\text{spike}} = U_{1:k} U_{1:k}^\top x, \qquad
x_{\text{tail}} = U_{k+1:d} U_{k+1:d}^\top x.
\]

The spike component lies in the dominant low-rank subspace and typically captures shared structural information such as syntax, formatting conventions, and global discourse patterns. The tail component lies in the residual subspace and encodes rarer, more semantically specific variations, such as entity-level distinctions or nuanced contextual cues.

Importantly, the spike component typically dominates in $\ell_2$ energy:

\[
\|x_{\text{spike}}\|_2^2 \gg \|x_{\text{tail}}\|_2^2,
\]

meaning that most variance of the representation is concentrated in the low-rank subspace. However, this fact alone does not determine which component is responsible for coordinate-level outliers.

\section{Empirical Observation: Outliers Are Tail-Dominated}

To determine which component drives activation outliers, we conduct an empirical analysis on model activations. For a given layer and training step, we collect activation matrices of shape $N \times d$ (for example, $8192 \times 1024$). After mean removal, we compute the top $1\%$ principal directions via PCA and decompose each activation into spike and tail components as described above.

We then define outliers as the top $0.1\%$ largest-magnitude entries in the activation matrix. For these outlier coordinates, we measure the relative contribution of spike and tail components. The results consistently show the following pattern.

At earlier training stages (e.g., step 10k), spike and tail components both contribute to extreme values. However, at later stages (e.g., step 180k), outliers are overwhelmingly dominated by the tail component. In some layers, more than $85\%$ of the squared magnitude of outlier entries arises from $x_{\text{tail}}$, while the spike contribution becomes negligible. Moreover, examining the single largest activation entry often reveals that nearly all of its magnitude originates from the tail component.

This leads to a crucial empirical conclusion:

\[
\|x_{\text{tail}}\|_\infty \gg \|x_{\text{spike}}\|_\infty,
\]

even though

\[
\|x_{\text{spike}}\|_2^2 \gg \|x_{\text{tail}}\|_2^2.
\]

In other words, the spike subspace dominates total variance, but the tail subspace dominates coordinate-level extremes.

\section{Implications for Low-Bit Quantization}

Under uniform symmetric $b$-bit quantization, the step size $\Delta$ is determined by the maximum absolute value:

\[
\Delta = \frac{\max_i |x_i|}{2^{b-1} - 1}.
\]

Since $\max_i |x_i|$ is largely determined by tail outliers, the quantization resolution is effectively governed by $x_{\text{tail}}$. This has two important consequences.

First, the quantization noise magnitude scales with $\Delta$, meaning that all coordinates, including spike components, suffer from coarser discretization due to tail outliers. Second, because $x_{\text{tail}}$ has smaller $\ell_2$ energy relative to $x_{\text{spike}}$, its signal-to-noise ratio (SNR) under quantization becomes much lower:

\[
\mathrm{SNR}_{\text{tail}} =
\frac{\|x_{\text{tail}}\|_2^2}{\|\epsilon\|_2^2}
\ll
\mathrm{SNR}_{\text{spike}}.
\]

Thus, low-bit quantization disproportionately harms the tail component, which contains rare semantic information. Structural spike information remains comparatively stable, while nuanced semantic distinctions degrade first.

\section{Activation Range Control as a Mitigation Strategy}

Given that tail-driven outliers inflate dynamic range and amplify quantization noise, it is natural to consider controlling activation magnitude during training. A principled method is Scheduled Soft Clipping (SSC), applied after LayerNorm:

\[
\tilde{z}_i = \alpha(t) \tanh\left(\frac{z_i}{\alpha(t)}\right).
\]

Here, $\alpha(t)$ is a time-dependent clipping radius. Early in training, $\alpha(t)$ is large, and the transformation behaves nearly linearly:

\[
\tilde{z} \approx z.
\]

Later in training, $\alpha(t)$ decreases, bounding activations and suppressing extreme tail excursions.

This approach has three desirable properties. First, it explicitly bounds coordinate magnitude, ensuring

\[
|\tilde{z}_i| \le \alpha(t).
\]

Second, it preserves gradient stability, since the derivative

\[
\frac{\partial \tilde{z}_i}{\partial z_i}
=
\mathrm{sech}^2\left(\frac{z_i}{\alpha(t)}\right)
\]

remains bounded between zero and one. Third, it mitigates tail-driven dynamic range inflation without destroying the low-rank spike structure.

\section{Conclusion}

The central insight of this chapter is that activation outliers in trained Transformer models are not primarily caused by dominant spike components but by long-tail subspace projections. Although the spike subspace carries most $\ell_2$ energy and encodes shared structural information, coordinate-level extremes are overwhelmingly driven by tail components, especially in later training stages. Because uniform quantization step size is determined by these extreme coordinates, low-bit training disproportionately harms tail semantics.

Activation range control through scheduled soft clipping provides a principled continuation-based strategy to suppress tail-driven outliers while preserving spectral structure. This approach improves quantization robustness and protects rare semantic features without collapsing the representation geometry. The interplay between spectral anisotropy, tail-dominated outliers, and quantization sensitivity provides a coherent theoretical foundation for robust low-bit Transformer training.